% azhyphen user manual. LPPL 1.3c or later; Anar Abdullayev.
% Build with make doc; the release builder inserts canonical model statistics.
\documentclass[11pt,a4paper]{article}
\usepackage{fontspec}
\usepackage[margin=25mm]{geometry}
\setmainfont{CMU Serif}
\setmonofont{CMU Typewriter Text}
\usepackage[hidelinks]{hyperref}
% Generated from VERSION; do not edit.
\def\azhyphenversion{0.2.0}
\endinput

% Keep PDF streams inspectable and suppress volatile PDF metadata.
\pdfvariable compresslevel=0
\pdfvariable objcompresslevel=0
\pdfvariable suppressoptionalinfo=767
\setlength{\parindent}{0pt}
\setlength{\parskip}{6pt}
\setlength{\emergencystretch}{2em}
\title{\textsf{azhyphen}\\Standalone Azerbaijani hyphenation}
\author{Anar Abdullayev}
\date{Package version \azhyphenversion}
\begin{document}
\maketitle

\section{Purpose and status}
\textsf{azhyphen} activates Azerbaijani hyphenation globally when loaded,
with the same document interface for LuaLaTeX and XeLaTeX. It uses each
engine's native hyphenation facilities and requires neither Babel nor
Polyglossia. Automatic language switching is outside its scope.

Versions below 1.0.0 are development releases. The package is usable, but
linguistic coverage is experimental and may change between releases.
Software versions are independent of pattern counts and evaluation scores.

\section{Requirements}
Use a working TeX Live installation with LuaLaTeX or XeLaTeX and the
\texttt{iftex} package. LuaLaTeX also requires \texttt{luaotfload}.
pdfLaTeX and other engines are unsupported and produce a package error.
XeLaTeX needs the named \texttt{azhyphen} language registered in its format
before a document is compiled. Select a font containing Azerbaijani letters;
the package does not select a document font.

The repository's automated installation requires a POSIX shell, GNU Make,
Python 3.10 or newer, and TeX Live's \texttt{kpsewhich}, \texttt{tlmgr}, and
\texttt{fmtutil-user}. Installation is tested with TeX Live 2026 on Ubuntu
and native Windows with MSYS2. Other TeX distributions have not been verified.

\section{Installation}
When available, install via your TeX distribution's package manager.
For a repository checkout, \texttt{make install} copies the runtime to your
user TeX tree, registers XeLaTeX, rebuilds its format, and verifies both
engines. Installation is repeatable and requires no administrator access.
\texttt{make uninstall} removes only package-owned files and registration.

\newpage
\begin{samepage}
\subsection{Installing the CTAN archive manually}
The CTAN archive contains public runtime and documentation files; it does
not contain the repository's Makefile or maintenance scripts. Obtain your
user tree with \texttt{kpsewhich -var-value=TEXMFHOME}. Copy these files
into its \texttt{tex/latex/azhyphen} directory:
\begin{quote}\ttfamily
azhyphen.sty, azhyphen.lua, azhyphen-version.tex,\\
azhyphen.pat.txt, azhyphen.hyp.txt, azhyphen-xetex.tex
\end{quote}
\end{samepage}
Keep \texttt{azhyphen-data.json} with them as a checksum manifest.
Store the manual and licensing files in the user tree's
\texttt{doc/latex/azhyphen} directory. LuaLaTeX can then load the patterns
directly. Refresh the filename database if your distribution requires it.

For XeLaTeX, preserve existing user language configuration. The user
tree's \path{tex/generic/config} directory may contain
\path{language-local.dat} and \path{language.dat}. Back up the existing
files before editing. Add exactly one line to \path{language-local.dat}:
\begin{verbatim}
azhyphen azhyphen-xetex.tex
\end{verbatim}
If the name \texttt{azhyphen} already names another loader, resolve the
conflict first. If a user \texttt{language.dat} exists, append the same
line there once, preserving all unrelated entries. Otherwise generate it
using TeX Live, substituting your user tree for \texttt{TREE}:
\begin{verbatim}
tlmgr generate language.dat \
  --localcfg TREE/tex/generic/config/language-local.dat \
  --dest TREE/tex/generic/config/language.dat
fmtutil-user --byfmt xelatex
\end{verbatim}
Run the format rebuild after either method. Missing unrelated language
loaders indicate a TeX installation problem; repair that installation and
retry. To undo a manual installation, remove only this registration line
from both configuration files and rebuild the user format; remove the
package files. Restore backups only if no unrelated edits occurred since.
The repository installer provides automated ownership and rollback handling
and is preferable when available.

\section{Basic usage}
\begin{verbatim}
\documentclass{article}
\usepackage{fontspec}
\setmainfont{CMU Serif} % an Azerbaijani-capable font
\usepackage{azhyphen}
\begin{document}
Azərbaycan dilində mətn avtomatik olaraq hecalanır.
\end{document}
\end{verbatim}
Compile with \texttt{lualatex} or \texttt{xelatex}. Load the package with
\verb|\usepackage{azhyphen}|. No package options or document language
configuration are required.

\subsection{LuaLaTeX}
At load time the Lua runtime locates the bundled canonical patterns via
kpathsea, creates a native LuaTeX language, installs its hyphenation
character mappings, loads any exceptions, and selects it globally.
A missing, empty, or invalid pattern file causes a clear load error.

\subsection{XeLaTeX}
The format loads the mechanically generated pattern loader under the
name \texttt{azhyphen}. Package loading selects that language globally
and sets Azerbaijani case mappings. In particular, dotted i/İ and
dotless ı/I remain distinct. If registration is absent, loading fails
with an installation diagnostic. Updating the patterns or exceptions
requires rebuilding the user XeLaTeX format.

\section{Hyphenation examples}
To inspect the current engine's permitted breakpoints in the compilation
log, add this to a document using the package:
\begin{verbatim}
\showhyphens{fotoelektrik ellərarası addımölçən}
\showhyphens{videobloger ətürpərdici}
\end{verbatim}
These are diagnostic sample words, not authoritative syllabification
claims. Actual line breaks depend on the surrounding text, font, and
line width. Review important typography manually.

Both left and right hyphenation minima are 2. These values are provisional
pending Azerbaijani linguistic and typographic review. For a document-local
reviewed correction, standard TeX \verb|\hyphenation{...}| exceptions or
explicit \verb|\-| discretionary breaks may be used after loading the
package. They do not change the bundled model.

\section{Bundled data and quality}
\texttt{azhyphen.pat.txt} is the single canonical PATGEN pattern dataset.
\texttt{azhyphen-xetex.tex} is mechanically derived from it, not a separate
linguistic source. \texttt{azhyphen.hyp.txt} is intentionally empty until
reviewed exceptions exist. The data manifest records checksums of the
canonical patterns and derived loader.

Public statistics come from the matching localization snapshot and are
inserted into this manual during the build. Evaluation uses held-out
manually annotated words that were not used for training.

% AZHYPHEN-DOC-STATS:START
\emph{Run make doc to insert the synchronized public model statistics.}
% AZHYPHEN-DOC-STATS:END

These metrics describe the bundled experimental model, not complete
coverage of Azerbaijani. Proper names, compounds, loanwords, and other
unseen words may need manual correction. The package supplies no language
switching, translation, or other language-management features.

\section{Provenance, license, and support}
The patterns and public statistics originate in Anar Abdullayev's separate
localization project. This distribution contains a validated snapshot;
runtime use requires no access to the training project. The exact source
commit and distribution grant are recorded in \texttt{PROVENANCE.md}
(\texttt{docs/PROVENANCE.md} in the repository). Training inputs and
evaluation archives are outside this package's distribution scope.

The maintained work, including bundled patterns and documentation, is
licensed under the LaTeX Project Public License 1.3c or later. See
\texttt{LICENSE} and the complete \texttt{LPPL-1.3c.txt}. The current
maintainer is Anar Abdullayev.

Source and bug reports: \url{https://github.com/abdanar/azhyphen}.
Report package or installation problems with the engine, TeX distribution,
and a minimal reproduction. Report linguistic issues separately, showing
the word, observed breakpoints, and the proposed correction with evidence.
Do not include private documents or corpus inputs.
\end{document}
